\documentclass{article}

\usepackage[english]{babel}
\usepackage[letterpaper,top=2cm,bottom=2cm,left=3cm,right=3cm]{geometry}
\usepackage{amsmath,amssymb}
\setlength{\parindent}{0pt}

\title{DSC 190/291: Learning Theory through Formal Proof and Proof Presentation\\Assignment 7}
\date{UCSD, Spring 2026}

\begin{document}
\maketitle

\section{Overview}

This assignment has five parts:
\begin{itemize}
\item Part A: kernel methods.
\item Part B: $\ell_1$ Rademacher complexity.
\item Part C: normalized $\ell_1$ margins and FixedBoost.
\item Part D: a FixedBoost margin experiment.
\item Part E: agent skill reflection.
\end{itemize}

\section{Part A: Kernel Methods}

Let $\mathcal{U}$ be a Hilbert space with inner product $\langle \cdot,\cdot \rangle$ and norm $\|\cdot\|$, and let
$$
\phi : \mathcal{X}\to \mathcal{U}
$$
be a feature map with kernel
$$
K(x,x')=\langle \phi(x),\phi(x')\rangle.
$$
For $B>0$, define
$$
\mathcal{H}_{K,B}
=
\left\{
x \mapsto \langle w,\phi(x)\rangle
:
w\in\mathcal{U},\ \|w\|\le B
\right\}.
$$
For a sample $S=(x_1,\dots,x_n)$ and real-valued function class $\mathcal{F}$, write
$$
\widehat{\mathfrak{R}}_S(\mathcal{F})
=
\mathbb{E}_\sigma
\sup_{f\in\mathcal{F}}
\frac{1}{n}\sum_{i=1}^n \sigma_i f(x_i),
$$
where $\sigma_1,\dots,\sigma_n$ are independent uniform signs, and
$$
\mathfrak{R}_{\mathcal{D},n}(\mathcal{F})
=
\mathbb{E}_{S\sim \mathcal{D}^n}\widehat{\mathfrak{R}}_S(\mathcal{F}).
$$

\begin{enumerate}
\item Prove
$$
\widehat{\mathfrak{R}}_S(\mathcal{H}_{K,B})
\le
\frac{B}{n}\sqrt{\sum_{i=1}^n K(x_i,x_i)},
$$
and conclude
$$
\mathfrak{R}_{\mathcal{D},n}(\mathcal{H}_{K,B})
\le
\frac{B\sqrt{\mathbb{E}[K(X,X)]}}{\sqrt{n}}.
$$

\item Let
$$
J(w)
=
\frac{1}{n}\sum_{i=1}^n \operatorname{loss}(\langle w,\phi(x_i)\rangle;y_i)
+
\lambda \|w\|^2
$$
with $\lambda>0$. Prove that any minimizer can be chosen in the form
$$
w^\star = \sum_{j=1}^n \alpha_j \phi(x_j),
\qquad
\alpha \in \mathbb{R}^n,
$$
and rewrite $J$ as an optimization problem over $\alpha$ using only the Gram matrix $K_S$.
\end{enumerate}

\section{Part B: \texorpdfstring{$\ell_1$}{l1} Rademacher Complexity}

\begin{enumerate}
\item For a finite class $\mathcal{F}$ of real-valued functions, prove
$$
\widehat{\mathfrak{R}}_S(\operatorname{conv}(\mathcal{F}))
=
\widehat{\mathfrak{R}}_S(\mathcal{F}).
$$

\item Let $\phi : \mathcal{X}\to \mathbb{R}^d$ satisfy
$$
\|\phi(x)\|_\infty \le R
$$
for all $x$, and define
$$
\mathcal{H}^1_B
=
\left\{
x \mapsto \langle w,\phi(x)\rangle
:
\|w\|_1 \le B
\right\}.
$$
Derive a logarithmic-in-$d$ Rademacher bound for $\mathcal{H}^1_B$.

\item Using the Week 7 uniform-convergence theorem, derive a uniform bound for
$$
L^\ell_{\mathcal{D}}(w)-L^\ell_S(w)
$$
over $\|w\|_1\le B$ for a $G$-Lipschitz loss $\ell$ bounded in $[0,1]$, and compare it with the Week 4 sparsity/VC bound.
\end{enumerate}

\section{Part C: \texorpdfstring{Normalized $\ell_1$ Margins and FixedBoost}{Normalized l1 Margins and FixedBoost}}

Let
$$
S=((x_1,y_1),\dots,(x_n,y_n)),
\qquad
y_i\in\{-1,+1\},
$$
and let $\mathcal{G}$ be a finite symmetric base class of binary predictors. For coefficients $w=(w_g)_{g\in\mathcal{G}}$, define
$$
f_w(x)=\sum_{g\in\mathcal{G}} w_g g(x),
\qquad
\operatorname{margin}_S(w)=\min_i \frac{y_i f_w(x_i)}{\|w\|_1}.
$$

FixedBoost uses a fixed step size $\eta>0$ and rounds $t=1,\dots,T$:
\begin{itemize}
\item choose $g_t$ maximizing the weighted correlation $\sum_i D_{t,i} y_i g_t(x_i)$,
\item update $w_t=w_{t-1}+\eta e_{g_t}$,
\item reweight examples by $D_{t+1,i}\propto \exp(-y_i f_{w_t}(x_i))$.
\end{itemize}

\begin{enumerate}
\item Prove a population $0/1$ error bound from an empirical normalized margin lower bound $\operatorname{margin}_S(w)\ge \gamma$.
\item Use the FixedBoost black-box theorem to derive a generalization guarantee for the classifier it returns.
\item Encode the sample/base class by the matrix
$$
A_{ig}=y_i g(x_i)\in\{-1,+1\}^{n\times |\mathcal{G}|}.
$$
Use von Neumann's minimax theorem to prove the exact relationship between the best weak-learning edge and the best normalized $\ell_1$ margin.
\item Explain how this point of view refines the Homework 6 sparse-boosting analysis.
\end{enumerate}

\section{Part D: FixedBoost Margin Experiment}

The base class is the $2d$ signed coordinate predictors on Boolean inputs $x\in\{-1,+1\}^d$.

\begin{enumerate}
\item Implement FixedBoost from scratch directly on the sign matrix $A$ with fixed step size.
\item Formulate and solve the optimal normalized-margin LP, and compare the FixedBoost trajectory with the optimum on at least two instances: one easy and one hard.
\item Interpret the convergence behavior and compare it with the theorem from Part C.2.
\end{enumerate}

\section{Part E: Agent Skill Reflection}

\begin{enumerate}
\item Describe which existing workflow or skill you used.
\item Describe whether you developed or modified a skill; if not, propose one.
\item Explain how you checked the math, code, plots, optimization outputs, and final claims.
\end{enumerate}

\end{document}
